\documentclass[11pt,a4paper]{article}
% When coordination pays — the consolidated write-up of the study. Figures are exported from the public site
% (site/scripts/export-figures.mjs) and tables generated from the run data (paper/tables.py); `make paper` does both.
% LuaLaTeX: Libertinus text and math cover the Greek and symbols used in the generated glossary.
\usepackage{amsmath}
\usepackage{libertinus}
\usepackage[margin=2.6cm]{geometry}
\usepackage{microtype}
\usepackage{graphicx}
\usepackage{booktabs}
\usepackage{xcolor}
\usepackage[font=small,labelfont=bf]{caption}
\usepackage{subcaption}
\usepackage{enumitem}
\usepackage[round,authoryear]{natbib}
\usepackage[colorlinks=true,linkcolor=accent,citecolor=accent,urlcolor=accent]{hyperref}
\usepackage[capitalise,noabbrev]{cleveref}

\definecolor{accent}{HTML}{8A5216}
\graphicspath{{figures/}}
\counterwithin{figure}{section}
\counterwithin{table}{section}
\setlist{itemsep=2pt,topsep=4pt}
\setlength{\parskip}{3pt}

\newcommand{\Thalf}{T_{1/2}}
\newcommand{\Teff}{T_{\mathrm{eff}}}
\newcommand{\tauE}{\tau_E}
\newcommand{\taum}{\tau_m}
\newcommand{\figw}[4]{\begin{figure}[tbp]\centering\includegraphics[width=#2]{#1}\caption{#3}\label{fig:#4}\end{figure}}
\newcommand{\figpair}[6]{\begin{figure}[tbp]\centering
  \begin{subfigure}{0.49\linewidth}\includegraphics[width=\linewidth]{#1}\caption{#2}\end{subfigure}\hfill
  \begin{subfigure}{0.49\linewidth}\includegraphics[width=\linewidth]{#3}\caption{#4}\end{subfigure}
  \caption{#5}\label{fig:#6}\end{figure}}
\newcommand{\takeaway}[1]{\par\smallskip\noindent\fcolorbox{accent!40}{accent!6}{\parbox{\dimexpr\linewidth-2\fboxsep-2\fboxrule}{\small #1}}\par\smallskip}

\title{\textbf{When coordination pays}\\[4pt]\large Attention networks under a moving, noisy source:\\regimes, costs, mechanisms and implications}
\author{Claude (Opus 5.5) \and Max Bromberg}
\date{\today}

\begin{document}
\maketitle
\begin{abstract}
\noindent
A group of agents learns whom to listen to while the source of useful news moves and the signal is noisy. Each agent
splits one unit of attention among the others and strengthens links that bring it the news soonest, with
self-reinforcement $\kappa$, commitment (selection temperature) $T$ and memory $\varepsilon$. We map the behaviour of
this model by simulation and solve it where possible. Coordination works in a window of commitment: below $T=\kappa$
the group locks onto whoever led initially, decided by a separatrix in the starting attention; above
$\Thalf=\kappa+(1-\kappa)\ln\rho/\ln(N-2)$ attention spreads; in between a single hub forms on the current source and
follows it. A hub costs about twice as much as no structure, measured by broadcast dissipation, at every group size,
while its benefit (log-score gain) grows like $\ln(N-1)$, so coordination has increasing returns to size. Noise is paid
in precision rather than structure: false alarms do not raise the cost but shrink the usable window by a similar fraction
at every size. Near the lock-in edge, tracking survives noise far better than mean-field theory predicts; switching
effects off one at a time shows that unequal rivals (a concavity effect) and relays through the hub account for it.
Rewarding a link for its position in the network beats rewarding direct evidence under noise, and comes within 1--3\%
of an ideal Bayesian observer up to 50\% false alarms; beyond that, the group's short memory is the limit. We relate
these results to the three lines of work the model joins: source-to-structure networks, optimal learning in noisy
niches, and adaptive transport networks.
\end{abstract}

\tableofcontents
\clearpage

\section{Introduction}\label{sec:intro}

Groups have to decide whom to listen to. A team follows whoever has the freshest information about a client; a
community of experts defers to whoever has been right lately; a flock follows the animals that seem to know where the
food is. In each case the source of useful information moves, and the signals about where it is are noisy. A group that
commits to a hub gains when the hub is right and pays when it is not; a group that never commits gains little but risks
little. This paper asks when such coordination pays, what it buys, what it costs, and what breaks it.

We study a minimal model. $N$ agents each split one unit of attention among the others. Every round a piece of news
appears at one agent, the source; with probability $\nu$ it appears at a random agent instead (a false alarm). Each agent
then strengthens the links that would have brought it the news soonest, counting paths through the network, with
self-reinforcement $\kappa$ (strong links earn more), commitment $T$ (how hard it bets on its best link) and memory
$\varepsilon$ (how fast old attention is replaced). The source stays put for $\tauE$ rounds on average and then moves.

The model joins three lines of work. The source-to-structure model of \citet{bromberg2020hierarchy} asked how the place
where information appears shapes who listens to whom, measuring closeness by a random-walk effective distance
\citep{brockmann2013hidden} and shape by hierarchy coordinates \citep{corominas2013originsheirarchy}. Optimal learning
in noisy niches \citep{lee2024constructing} asked how much memory a learner should keep when its environment switches,
and what happens when learners change the environment they track. Adaptive transport networks
\citep{corson2010fluctuations,katifori2010damage,ronellenfitsch2016global} grow the links that carry flow under a
budget, and form trees or loops depending on how the load fluctuates. No familiarity with any of them is assumed; the
model is defined from scratch in \cref{sec:model}, and \cref{sec:origins} says how each original premise fared.

\paragraph{Contributions.}
\begin{enumerate}[label=(\roman*)]
  \item A computed regime map with exact boundaries: lock-in for $T<\kappa$, coordinated hierarchy for
    $\kappa<T<\Thalf$, spread attention above $\Thalf$ (\cref{sec:regimes}). The one-variable reduction behind the
    boundaries is exact at any walk cost (A1), it remains the linear stability of the full network despite a kink at the
    fixed point (A2), and lock-in below $\kappa$ is decided by a separatrix in the initial weights (A3).
  \item Benefit and cost of coordination as functions of group size and noise (\cref{sec:scaling}): structural cost is
    flat in $N$ while benefit grows like $\ln(N-1)$; noise narrows the usable window of commitment instead of raising the
    cost.
  \item A mechanism for noise robustness near the lock-in edge, isolated by a factorial test (\cref{sec:robustness}):
    unequal rivals compete less because the score is concave in weight, and the hub relays even false news.
  \item A comparison of reward rules (\cref{sec:learning}): position (network closeness) against direct evidence
    (likelihood) and an ideal Bayesian observer.
  \item Results on memory, environmental pace and paid attention (\cref{sec:memory}), and on budgets and network shape
    against transport networks (\cref{sec:structure}).
\end{enumerate}

A companion website presents the same material as interactive chapters with enlargeable figures. All simulations,
analyses and checks are reproducible from the repository (\cref{app:methods}).

\section{The model}\label{sec:model}

\subsection{Agents and attention}
There are $N$ agents. The state of the group is an attention matrix $A$ in which $A_{ki}\ge 0$ is the share of agent
$i$'s attention that goes to agent $k$; every column sums to one and $A_{ii}=0$. Read as a network, $A$ has a link
$k\to i$ of strength $A_{ki}$: information that reaches $k$ travels on to those who listen to $k$. A group in which
everyone listens mostly to one agent is a star with that agent as its hub; one in which everyone listens a little to
everyone is a complete graph with thin links.

\subsection{One round}
Every round $t$, all agents update at once in four steps (\cref{fig:round}).

\figw{round}{\linewidth}{One round of the model. Bars are one agent's view of the five agents it can listen to; $s$ is
the one closest to the news. The same steps repeat every round, for every agent at once.}{round}

\paragraph{1. News appears} at the seeded agent $s_t$: the true source with probability $1-\nu$, otherwise a uniformly
random agent (a false alarm). Agents do not know which.

\paragraph{2. Score the links.} Closeness to the news is the random-walk effective distance \citep{brockmann2013hidden}.
Row-normalize the attention so that $P_{sk}=A_{sk}/\sum_j A_{sj}$ is the share of $s$'s audience made up by $k$, the
probability that a walker carrying news from $s$ steps to $k$. With a cost $\delta$ per step,
\begin{equation}
  Z=(I-e^{-\delta}P)^{-1},\qquad D_{sk}=-\ln\frac{Z_{sk}}{Z_{kk}} .
\end{equation}
$D_{sk}$ is short when many likely walks lead from $s$ to $k$. For $\delta\gg 1$ (we use $\delta=10$) only the direct
step counts, $D_{sk}\approx\delta-\ln P_{sk}$, a surprisal along the network. The source's own distance is set to the
smallest relay distance divided by $\rho=2.6$, $D_{ss}=\min_{k\ne s}D_{sk}/\rho$, so first-hand news beats any relay. The
info score of the link $k\to i$ combines closeness and existing strength:
\begin{equation}
  S_{ki}=D_{s_t k}^{\,\kappa-1}\,A_{ki}^{\,\kappa}.
\end{equation}
With self-reinforcement $\kappa>0$, strong links score more.

\paragraph{3. Commit.} Scores become rewards by tempered selection,
\begin{equation}
  r_{ki}=\frac{S_{ki}^{1/T}}{\sum_j S_{ji}^{1/T}} .
\end{equation}
Low commitment temperature $T$ sends almost all reward to the best link; high $T$ spreads it. The original model used a
hard cut (only the top fraction $1-\varphi$ of each agent's links is rewarded), which tempered selection contains as
$T\to 0$.

\paragraph{4. Blend with memory.}
\begin{equation}
  A_{ki}\leftarrow\frac{A_{ki}+\varepsilon\,r_{ki}}{1+\varepsilon},
\end{equation}
an exponential average of rewards with memory $\taum=1/\ln(1+\varepsilon)$; each column remains a distribution. One
step is exactly a forward-Euler step of the logit flow $\dot A=R(A)-A$ with step $\varepsilon/(1+\varepsilon)$, and the
fixed points are column softmaxes of $-\ln D$ at the effective temperature $\Teff=(T-\kappa)/(1-\kappa)$, the same family
as the learning dynamics of \citet{sato2003coupled} (bridges B2, B3; \cref{app:derivations}).

\subsection{The environment}
The true source stays put for $\tauE$ rounds on average (each round it jumps to a uniformly random other agent with
probability $1/\tauE$). Agents only see the seeded agent, which is a false alarm with probability $\nu$
(\cref{fig:environment}). This is the switching, noisily observed environment of \citet{lee2024constructing}.

\figw{environment}{\linewidth}{The environment over 72 rounds. Top: the true source, which stays put for $\tauE$ rounds
on average. Bottom: where the news lands; outlined cells are false alarms (here $\nu\approx 0.3$).}{environment}

\subsection{Parameters and implementation}
\Cref{tab:params} lists the parameters. Rewards go only to existing links, each agent ranks its own candidates, and
$\varepsilon$ is free. The implementation was checked step by step against the original code (\cref{app:validation}).

\begin{table}[tbp]\centering\small
\caption{Parameters.}\label{tab:params}
\begin{tabular}{lll}\toprule
parameter & meaning & values used\\\midrule
$T$ & commitment (selection temperature); low $T$ bets hard & $0.05$--$2$\\
$\kappa$ & self-reinforcement: strong links earn more & $0$--$0.6$; $0.2$ in the cost studies\\
$\varepsilon$ & memory, $\taum=1/\ln(1+\varepsilon)$ & $1$ ($\taum\approx1.44$) unless varied\\
$\rho$ & advantage of first-hand news over the best relay & $2.6$\\
$\delta$ & cost per walk step in the effective distance & $10$\\
$\nu$ & probability of a false alarm & $0$--$0.75$\\
$\tauE$ & mean rounds between moves of the true source & $10$--$3000$; $300$ in the cost studies\\
$N$ & number of agents & $10$--$80$\\\bottomrule
\end{tabular}\end{table}

\section{Measuring coordination}\label{sec:measures}

\paragraph{Tracking and commitment.} Tracking is the mean share of attention the other agents put on the current true
source, averaged over the second half of a run: $1$ when everyone listens only to the source, $1/(N-1)$ when attention
is spread evenly. Commitment is the mean over agents of their largest single share, whoever it goes to. High commitment
with low tracking is the signature of confident betting on the wrong agent.

\paragraph{Benefit: attention as a bet.} Each agent's attention column is a bet on where the next news will appear.
When it appears at $s$, agent $i$ is charged the log score $-\ln A_{si}$; this log-loss is $\mathrm{KL}(\delta_s\,\|\,A_{\cdot
i})$ for a point mass on $s$. Uniform attention pays $\ln(N-1)$. The benefit is
\begin{equation}
  B=\ln(N-1)-\Big\langle -\ln A_{s_t^{\mathrm{true}} i}\Big\rangle_{i\ne s,\,t},
\end{equation}
at most $\ln(N-1)$ and strongly negative for confident bets on the wrong agent. It is the Kelly regret of the attention
bet and the quantity used by \citet{lee2024constructing}.

\paragraph{Information cost and calibration.} The information cost
$I=\langle\mathrm{KL}(A_{\cdot i}\,\|\,\mathrm{uniform})\rangle_i=\ln(N-1)-\langle H(A_{\cdot i})\rangle$ measures how much
an agent committed, not whether it was right: in the example of \cref{fig:cost-benefit} the right and the wrong bet both
cost $1.13$ nats. For a calibrated bettor the expected log-loss equals the entropy, so the calibration gap $B-I$ is near
zero; confident and wrong drives it far below zero.

\paragraph{Structural cost: attention as pipes.} Read the mutual attention $C_{ij}=(A_{ij}+A_{ji})/2$, scaled to sum to
one, as conductances. When news appears at $s$ it must reach everyone: inject one unit at $s$, withdraw $1/(N-1)$ at
every other agent, and let Kirchhoff's laws decide the currents $I_e$. The broadcast dissipation $\sum_e I_e^2/C_e$,
averaged over sources and divided by the same quantity for uniform attention, is the structural cost: $1$ for no
structure, about $2$ for a strong hub. A hub makes news from the hub cheap (spokes in parallel) but news from anyone else
expensive, since it must squeeze through one spoke. The measure is borrowed from adaptive transport networks, where it is
the quantity being minimized \citep{corson2010fluctuations}; the attention model never pays it, and we use it to price
what the model builds.

\figw{cost-benefit}{\linewidth}{Benefit and structural cost on a 7-agent example. Top: one agent's attention over the
other six ($s$ current source, $o$ old source) and what each bet earns. Bottom: the attention read as pipes (grey width
= conductance) carrying a broadcast from $s$ (coloured width = current). The hub holds $0.85$ of every other agent's
attention.}{cost-benefit}

\paragraph{Why not distance.} The original model judged networks by mean walk distance. Distance from $s$ depends on how
$s$'s audience is split, not on how much each listener gives to $s$: a network where everyone listens only to $s$ and one
where everyone listens a little to everyone have the same mean distance from $s$ (13.09 for every setting in E1). The
all-pairs version behaves erratically under noise (2.4, 6.0, 1.9 and 0.1 for $\nu=0$ to $0.75$ at $N=10$). Log-loss sees
concentration and is used throughout.

\paragraph{Shape: hierarchy coordinates.} Treeness $T_h\in[-1,1]$ (do paths converge on few top nodes?),
feedforwardness $F\in[0,1]$ (how much path weight avoids cycles?) and orderability $O\in[0,1]$ (how many nodes lie
outside cycles?) place a network in the hierarchy morphospace \citep{corominas2013originsheirarchy}. Weighted networks
are thresholded at 8 levels and the coordinates averaged, so a perfect star reads $7/8=0.88$. We compute treeness with
matrix powers in the path counts (eq.~14 of the original), correcting two errors in the code used by the 2020 thesis
(\cref{app:errata}). \Cref{fig:shapes} shows typical values.

\figw{shapes}{\linewidth}{Three shapes and their hierarchy coordinates (typical runs).}{shapes}

\section{Three regimes}\label{sec:regimes}

Along the commitment axis a group falls into one of three regimes (\cref{fig:regimes}):
\begin{itemize}
  \item \textbf{Lock-in} ($T<\kappa$): whoever is ahead gets further ahead regardless of where the news is. The group
    builds a hub and pays its full structural cost, but the hub sits on whoever led early; the benefit is far below
    zero. Under a moving source this is over-betting.
  \item \textbf{Coordinated hierarchy} ($\kappa<T<\Thalf$): one hub forms on the current source and follows it when it
    moves.
  \item \textbf{Spread out} ($T>\Thalf$): no one commits; little is paid and little is gained.
\end{itemize}

\figw{regimes}{\linewidth}{The three regimes, measured at $N=40$, $\tauE=300$, $\nu=0.25$, $\kappa=0.2$ and $T=0.15$,
$0.25$, $0.8$. Benefit as a share of the maximum; cost relative to uniform attention.}{regimes}

\subsection{The regime map}
Experiment E5 runs 30 agents over 14 commitments, four self-reinforcements and five source schedules (constant; Markov
with $\tauE=10,100,1000$; a new random source every round), 3000 rounds and three repeats each, and classifies each run
by commitment and tracking. \Cref{fig:e5} shows the map; the lines are the exact boundaries of \cref{sec:boundaries}, not
fits. Lock-in on stale hubs needs both $\kappa>0$ and a moving source.

\begin{figure}[tbp]\centering
\includegraphics[width=\linewidth]{e5-constant}\\[2pt]
\includegraphics[width=\linewidth]{e5-markov100}\\[2pt]
\includegraphics[width=\linewidth]{e5-random}
\caption{Regime map (E5) for a constant source (top), a source moving every 100 rounds on average (middle) and a new
random source every round (bottom). Solid line $T=\kappa$, dashed $\Thalf$.}\label{fig:e5}
\end{figure}

\subsection{Exact boundaries}\label{sec:boundaries}
In the symmetric state every agent puts the share $x$ on the source and $p=(1-x)/(N-2)$ on each other agent. The model
then reduces to one variable. The source scores $(D/\rho)^{\kappa-1}x^\kappa$ and each peer $D^{\kappa-1}p^\kappa$; the
common factor $D^{\kappa-1}$ cancels from the reward ratio, which is $(\rho^{1-\kappa}u^\kappa)^{1/T}$ with $u=x/p$.
Its fixed point and the half-attention temperature are
\begin{equation}
  u^*=\rho^{(1-\kappa)/(T-\kappa)},\qquad x^*=\frac{u^*}{u^*+N-2},\qquad
  \Thalf=\kappa+\frac{(1-\kappa)\ln\rho}{\ln(N-2)},
\end{equation}
with contraction $(1+\varepsilon\kappa/T)/(1+\varepsilon)$ per step, below one iff $T>\kappa$. The window
$(\kappa,\Thalf)$ has width $(1-\kappa)\ln\rho/\ln(N-2)$ and narrows as $N$ grows. Rescaled by it, the share on the
source for every group size falls on one curve (E2, \cref{fig:e2}).

The reduction is exact at any walk cost $\delta$ (A1): in the symmetric state the source's own column is uniform, so all
relays are equivalent under permutations fixing $s$. The update has a kink at its fixed point, because all relay
distances tie in $\min_k D_{sk}$; linearizing one smooth piece at a time, the full network's spectral radius equals the
reduced contraction up to a relay group 0.01--4\% above it, and crosses one at $T=\kappa$ (A2, \cref{tab:spectra}). No
instability hides in the full network.

\figw{e2-collapse}{0.8\linewidth}{Share of attention on a constant source against $z=\Teff\ln(N-2)/\ln\rho$
($\kappa=0$). Every size crosses one half at $z=1$ (E2).}{e2}

\begin{table}[tbp]\centering\small
\caption{Linearization of the anchored piece at the symmetric state, $N=10$, $\kappa=0.3$, $\varepsilon=1$ (A2). For
$T<\kappa$ the unstable interior point is used.}\label{tab:spectra}
\input{generated/spectra}
\end{table}

\subsection{Lock-in is decided at the start}\label{sec:separatrix}
Below $T=\kappa$ the symmetric state is unstable. For a column dominated by the source $s$ and the strongest rival $q$,
the log-ratio $l=\ln(w_s/w_q)$ has the fixed point $l^*=(1-\kappa)\ln\rho/(T-\kappa)$, unstable for $T<\kappa$. A listener
therefore locks onto its rival iff initially
\begin{equation}
  \frac{w_q}{w_s}>\rho^{(1-\kappa)/(\kappa-T)},
\end{equation}
which recovers the one-step criterion of the hard cut as $T\to0$ (A3). Listener by listener, this predicts the simulated
winner 95--100\% of the time (\cref{tab:separatrix}). The signal never gets a say: this is lock-in by historical accident
\citep{arthur1989competing} in network form.

\begin{table}[tbp]\centering\small
\caption{Separatrix prediction against simulation, $N=20$, constant source, 3 repeats (A3).}\label{tab:separatrix}
\input{generated/separatrix}
\end{table}

\takeaway{The safe setting is just above the lock-in edge: close enough to $\kappa$ to commit, far enough not to freeze.}

\section{Size and noise}\label{sec:scaling}

Experiment E8 runs groups of $N=10,20,40,80$ agents with a source moving every $\tauE=300$ rounds on average, false-alarm
rates $\nu=0,0.25,0.5,0.75$, seven commitments $T\in\{0.15,\dots,0.8\}$, $\kappa=0.2$ and $\varepsilon=1$, for 3000 rounds
and two repeats each. For every $(N,\nu)$ we take the commitment with the highest benefit (\cref{tab:e8}).

\figw{scaling}{\linewidth}{Benefit (nats) and structural cost against group size at the best commitment, for three noise
levels. The dashed line is the maximum benefit $\ln(N-1)$.}{scaling}

\paragraph{Increasing returns to size.} With the best commitment, coordination recovers 94--99\% of the possible benefit
up to $\nu=0.5$ at every size, and 67--75\% at $\nu=0.75$. Structural cost stays at 1.8--2.0 times that of uniform
attention at every size, so benefit per unit of cost rises from 1.2 to 2.3 nats between 10 and 80 agents
(\cref{fig:scaling,fig:e8pair}). Heavy noise lowers the cost to 1.4--1.7 because the group commits less; at 80 agents the
return per unit of cost is unchanged.

\figpair{e8-benefit}{Benefit as a share of $\ln(N-1)$.}{e8-per-cost}{Benefit per unit of cost.}{E8 at the best
commitment.}{e8pair}

\paragraph{Noise lowers what coordination buys.} Over the whole range of commitment, benefit against cost traces a
frontier; noise moves it down, not to the right (\cref{fig:frontier}). Noise is paid in commitment instead: the best
$T$ moves from $0.30$ to $0.25$, toward the lock-in edge, as agents must bet harder on the current hub to average out
false alarms --- the attention counterpart of the longer memory noise selects (\cref{sec:memory}).

\figpair{e8-frontier-10}{10 agents}{e8-frontier-80}{80 agents}{Structural cost against benefit over $T>\kappa$, one curve
per noise level.}{frontier}

\paragraph{The window narrows.} Tracking falls below one half at an effective $\Thalf(\nu)$ that decreases with noise
(\cref{tab:thalf}). The usable window $(\kappa,\Thalf(\nu))$ shrinks with group size, since $\Thalf\to\kappa$ like
$1/\ln N$, and with noise, by a similar fraction at every size: about 20\% at $\nu=0.25$, half at $0.5$ and 80\% at
$0.75$ (\cref{fig:window}). At 80 agents and $\nu=0.5$ commitment must lie within about $0.09$ of $\kappa$.

\figw{window}{\linewidth}{Usable range of commitment, from $T=\kappa$ to where tracking falls below one half, per group
size and noise level (lightest: $\nu=0$). At $\nu=0.75$ for 40 and 80 agents the end comes from the relay-aware walk,
since the runs never reached one half.}{window}

\figpair{e8-tracking-10}{10 agents}{e8-tracking-80}{80 agents}{Tracking against commitment. Near $\kappa$ tracking
survives noise up to $\nu=0.5$; toward $\Thalf$ it falls early.}{tracking}

\paragraph{Lock-in pays the full price for nothing.} At $T=0.15<\kappa$ the structural cost is the full 1.9--2.0 at every
noise level, but the hubs are stale and the benefit far below zero: the most expensive structure with the worst return.

\begin{table}[tbp]\centering\small
\caption{E8 at the best commitment for each group size and noise level.}\label{tab:e8}
\input{generated/e8}
\end{table}

\begin{table}[tbp]\centering\small
\caption{Where tracking falls to one half: closed form without noise, simulation (interpolated), relay-aware walk and
averaged mean field.}\label{tab:thalf}
\input{generated/thalf}
\end{table}

\takeaway{Large, noisy groups gain the most from coordinating and have the least room for error.}

\section{Why the hub survives noise}\label{sec:robustness}

\subsection{The puzzle}
Near the lock-in edge, tracking stays at 0.91--0.99 with up to 50\% false alarms at every size. A noisy mean field
averages the reward on the source over the noise in the symmetric state (A4):
\begin{equation}
  \mathbb E[r](x)=\Big(1-\nu+\tfrac{\nu}{N}\Big)f(x)+\nu\,\tfrac{N-1}{N}\,g(x),
\end{equation}
where $f$ is the reward when the true source is seeded and $g$ when a rival is (the rival then carries the $\rho$
advantage). The averaged map $x\mapsto(x+\varepsilon\mathbb E[r])/(1+\varepsilon)$ matches the simulation within 3\%
for $T\ge0.45$ but falls far short near $\kappa$ (\cref{tab:meanfield}). Keeping the round-to-round fluctuations (a
Jensen effect) and the relay effect closes part of the gap.

\begin{table}[tbp]\centering\small
\caption{Tracking: simulation (E8) against the averaged mean field, the fluctuating walk and the relay-aware walk.}\label{tab:meanfield}
\input{generated/meanfield}
\end{table}

\subsection{Switching effects off one at a time}
Every reduction keeps the symmetric state, in which all rivals of the source hold equal attention; the simulation does
not. A $2\times2$ factorial removes that one difference: after every round, each agent's attention on its rivals is reset
to their mean, leaving the share on the source unchanged. A second switch removes relays, by rewarding direct evidence
instead of closeness (\cref{sec:learning}). Against these runs we place four reductions: the symmetric walk, the
relay-aware walk, and a column walk that keeps one agent's full unequal column with the rest of the network symmetric
(\cref{tab:a4,fig:a4}).

\figw{mechanism}{\linewidth}{Why the hub keeps its place. Left: uneven rivals pull less in total
(score $\propto w^{\kappa/T}$, concave). Middle: news from a false alarm passes through the hub. Right: tracking at
$N=10$, $\nu=0.5$, $T=0.25$, removing one effect at a time.}{mechanism}

\figpair{a4-10}{10 agents}{a4-80}{80 agents}{The A4 test at $\nu=0.5$: tracking against $T$ for each condition and
reduction.}{a4}

\begin{table}[tbp]\centering\small
\caption{The A4 test at $T=0.25$. ``Equal'': rivals equalized every round. ``Competition'': total rival pull relative to
equal rivals in the evidence column walk.}\label{tab:a4}
\resizebox{\linewidth}{!}{\input{generated/a4}}
\end{table}

\paragraph{Results.}
\begin{enumerate}
  \item With rivals equalized, each simulation matches its reduction: within 0.03 at $N\ge20$, 0.06 at $N=10$.
  \item Unequal rivals carry the whole excess near the edge: at $N=10$, $\nu=0.5$, $T=0.25$ tracking falls from $0.94$ to
    $0.75$ when rivals are equalized; at $N=80$ from $0.94$ to $0.82$. Far from the edge nothing changes.
  \item Relays amplify the effect: unequal rivals add $0.20$ with relays and $0.09$ without.
\end{enumerate}

\subsection{Mechanism: concavity}
A link's pull grows like $w^{\kappa/T}$. Above the lock-in edge $\kappa/T<1$ and the pull is concave in weight. Each false
alarm deposits extra attention on one random rival, which decays over the next few rounds; the leftover attention is
therefore uneven, and uneven weight has less total pull than the same weight spread evenly. The competition factor
$\sum_k w_k^{\kappa/T}/[(N-2)\bar w^{\kappa/T}]$ is 0.87--0.94 at $T=0.25$ and approaches one as $T$ grows --- exactly
where the excess appears and disappears. A single unequal column recovers 40--100\% of the excess (typically two
thirds); the remainder comes from relay distances, which also become unequal.

\paragraph{Relays.} When a false alarm fires at $j$, closeness is measured from $j$. Everyone listens mostly to the hub,
so the hub is the closest relay even to $j$'s news and keeps more reward than an ordinary rival would.

\paragraph{The visible trace.} With noise, treeness holds at the plateau while feedforwardness falls (0.88 to
0.57--0.78 at $T=0.25$): the short-lived side hubs form loops among the rivals (\cref{fig:ff}).

\figw{e8-ff}{0.8\linewidth}{Feedforwardness against commitment, 80 agents. Noise lowers it even where tracking is still
near one.}{ff}

\takeaway{The loops that appear under noise are how the group absorbs false alarms; equalize them away and the hub
becomes fragile.}

\section{Learning through others against learning alone}\label{sec:learning}

So far each agent rewards a link for its \emph{position}: closeness to the news, counting paths through others. The
alternative rewards direct \emph{evidence}: the likelihood of the observation if the sender were the source. Experiment
E9 keeps the explicit factor $A^\kappa$ and replaces the distance factor by $\lambda^{1-\kappa}$ for the seeded agent and
1 for everyone else,
\begin{equation}
  S^{\mathrm{ev}}_{ki}=\lambda_k^{1-\kappa}A_{ki}^\kappa,\qquad \lambda_k=\begin{cases}\lambda & k=s_t\\ 1&\text{otherwise,}\end{cases}
\end{equation}
a Kelly or multiplicative-weights reward. With $\lambda=\rho$ the two rules are identical in the symmetric state (A1), so
any difference comes from relays. A second variant uses the likelihood ratio of the true noise model,
$\lambda^*=1+N(1-\nu)/\nu$. As a bound, an ideal statistician runs the exact forward (Bayes) filter over the source,
knowing $\nu$ and $\tauE$. Grid, seeds and environments are those of E8 (\cref{fig:rules,tab:e9}).

\figw{rules}{\linewidth}{Position against evidence after a false alarm (highlighted ring: who gets the credit), and benefit
at the best commitment for 80 agents.}{rules}

\figpair{e9-benefit-10}{10 agents}{e9-benefit-80}{80 agents}{Benefit against noise at the best commitment, as a share of
$\ln(N-1)$.}{e9}

\paragraph{Position wins under noise.} At $\nu=0.5$, position recovers 94--98\% of the possible benefit and evidence
86--91\%; tracking is 0.91--0.95 against 0.72--0.79. At $\nu=0.25$ they are within one point, and without noise
identical. Position comes within 1--3\% of the ideal statistician up to $\nu=0.5$. Structural cost is set by the structure,
not the rule: the ideal statistician builds the same hub at the same 1.8--2.0.

\paragraph{Tuning evidence to the noise backfires} (76--82\% at $\nu=0.5$). The update averages rewards; it does not
multiply beliefs as Bayes' rule does. A sharper likelihood therefore does not accumulate evidence; it makes every round
copy the latest signal.

\paragraph{Memory is where the ideal pulls ahead.} At $\nu=0.75$ position recovers 67--75\% against 90--97\% for the
statistician. With $\varepsilon=1$ the group averages over about 1.4 rounds; the statistician uses every signal since the
source last moved. The calibration gap (\cref{fig:calibration}) is near zero for the statistician and for the model's
rules at their best $T$, and about $-690$ nats below $\kappa$ (over-betting).

\figpair{e9-tracking}{Tracking, $N=40$, $\nu=0.5$.}{e9-calibration}{Calibration gap $B-I$, $N=40$, $\nu=0.5$.}{Reward rules
against commitment.}{calibration}

\begin{table}[tbp]\centering\small
\caption{Each reward rule at its best commitment (E9).}\label{tab:e9}
\input{generated/e9}
\end{table}

\paragraph{What each rule stands for.} Position rewards being close to where news appears, directly or through others:
routing and social learning, as in citation, rumour or supply networks, with collective error correction through relays.
Evidence rewards the observation's likelihood: individual inference, as by Kelly bettors, experts' weights or Bayesian
learners, and as by the agents of \citet{lee2024constructing}.

\takeaway{Learning through intermediaries is cheap error correction when memory is short; where memory can be long,
careful individual inference is the benchmark.}

\section{Memory, pace and payment}\label{sec:memory}

\subsection{Memory should grow with the pace of change}
A learner in a changing world trades responsiveness against noise: a short memory follows changes but is thrown about by
noise; a long memory averages noise out but lags behind changes. For a single learner the optimal memory grows like
$\tauE^{1/2}$ \citep{lee2024constructing}. The attention network has two candidate memories: the blending weight
$\varepsilon$, with $\taum=1/\ln(1+\varepsilon)$, and self-reinforcement $\kappa$, which also makes the past persist. With a
perfect signal the shortest memory always wins, so E1 adds $\nu=0.5$ and scores runs by log-loss on the true source
(30 agents, $T=0.3$, 15 memories, six paces, four repeats). With $\varepsilon$ as the memory, $\taum^*\propto\tauE^{0.43}$,
approaching $1/2$ from below; with $\kappa$ it barely moves (\cref{fig:e1,tab:e1}). Self-reinforcement is a commitment
knob, not a memory --- which is why it sets the lock-in edge.

\figw{e1}{0.8\linewidth}{Best memory against the pace of change (E1), logarithmic axes.}{e1}

\begin{table}[tbp]\centering\small\caption{E1: optimal memory per pace.}\label{tab:e1}\input{generated/e1}\end{table}

\subsection{When memory sets the pace}
In niche-learning models agents can also change their environment: stabilizers slow the switching they track,
destabilizers speed it up \citep{lee2024constructing}. If a population's memory sets how fast its own world changes,
does the optimal memory settle, cycle or drift? E4 uses the binary environment of \citet{lee2024constructing} with a
mean-field population: for each resident memory, rare invaders of every other memory are tested, and the best response
is iterated (\cref{fig:e4,tab:e4}). In every case iteration converges to a stable memory from both ends, with no cycles or
drift. At that point the memory equals a passive agent's optimum for the environment the population creates, so the
square-root scaling closes on itself. Stabilizers slow their world 15--18-fold. The stable memory is shorter than the one
best for the group (10.7 against about 14 at $\tauE=1000$): shorter-memory invaders free-ride on the calm residents
provide, a quantitative version of the commons argument of \citet{lee2024constructing}.

\figw{e4}{0.8\linewidth}{Best response against the residents' memory, $\tauE=1000$ (E4). Crossings of the dashed identity
are self-consistent memories; all are stable.}{e4}

\begin{table}[tbp]\centering\small\caption{E4: self-consistent memories.}\label{tab:e4}\input{generated/e4}\end{table}

\subsection{Pooling signals}
Agents could pool several rounds of news before each update. Summing distances over the window never helps, because
the source's advantage is the $\rho$ contrast, which summing dilutes. Averaging rewards over a window of $k$ signals does
help: it acts as a sampling window whose optimum grows like $\tauE^{0.42}$, close to the memory scaling (E7,
\cref{fig:e7}).

\figw{e7}{0.8\linewidth}{Log-loss against the pooling window at $\tauE=300$ (E7).}{e7}

\subsection{Paying for attention does not entrench hubs}
If being listened to pays, do hubs use it to stay hubs, as monopoly arguments suggest? In E3 agents earn in proportion to
the attention they relay, and each epoch the worst earner copies the best earner's $\kappa$. Hubs do not entrench: they
change every one or two epochs with or without payment; with a slow source, payment if anything shortens their tenure
(\cref{tab:e3}). The same holds when the next source is drawn in proportion to relayed attention, and on sparse graphs.
Unlike a stabilizer, a hub cannot make itself more reliable by holding attention: nothing in the model lets it change the
source's pace.

\begin{table}[tbp]\centering\small\caption{E3: hub tenure and takeovers over 150 epochs.}\label{tab:e3}\input{generated/e3}\end{table}

\section{Budgets and shapes}\label{sec:structure}

\subsection{Attention against transport}
Adaptive transport networks also strengthen links that are used \citep{tero2010rules}. Under a steady load they form
trees; fluctuating loads keep loops \citep{katifori2010damage,corson2010fluctuations}. E6 writes both as one update,
$W\leftarrow(1-\eta)W+\eta\,\mathcal N[U^p]$, and swaps two ingredients (\cref{fig:e6flow}): the budget $\mathcal N$ (per
agent, as attention, or one shared pool, as material) and the usage $U$ (the info score, or the squared Kirchhoff current
when the source broadcasts). The per-agent/closeness cell is the attention model with tempered selection ($T=1/p$); the
shared/flow cell has Corson's optimum as its steady-load fixed point ($\gamma=1/p-1$). Runs use a $7\times7$ lattice,
$p\in[0.4,20]$, $\eta=0.05$, 3000 steps, three repeats, with constant, Markov and random sources.

\figw{e6-flow}{\linewidth}{The shared update of E6; the four cells are the four choices of usage (box 2) and budget
(box 4).}{e6flow}

\figpair{e6-loops-constant}{constant source}{e6-loops-random}{random source}{Independent loops of the backbone against
concentration $p$ (maximum 36). Shared-budget points that collapse onto part of the graph are omitted.}{e6loops}

\begin{itemize}
  \item \textbf{Trees need self-reinforcement.} The flow reward reinforces itself; closeness does so only through
    $\kappa$. At $\kappa=0$ equally close parents tie forever and every loop survives; trees appear past $p\approx1/\kappa$,
    where the symmetric state loses stability ($T=\kappa$).
  \item \textbf{A fluctuating source keeps loops only with the flow reward.} With closeness, a random source merely delays
    the transition; past it the network locks into some tree whatever the source does.
  \item \textbf{The budget decides robustness and cost.} A per-agent budget forces every agent to keep a link: strong
    concentration gives trees, never a collapse, with dissipation at most about 4.5 times the uniform lattice. A shared
    budget puts the tree transition exactly at Corson's $\gamma=1$ but collapses at strong concentration, with dissipation
    growing by orders of magnitude.
\end{itemize}

\subsection{Paths in the hierarchy morphospace}
Following one experimental variable at a time through $(T_h,F,O)$ shows where each experiment moves the shape
(\cref{fig:paths}). Commitment is a one-dimensional descent from the hierarchy corner to about $(0.25,0.25,0.25)$ as $T$
passes $\Thalf$; noise lowers feedforwardness first; and a hard selection cut climbs toward the same region under a random
source as under a constant one, so that structure comes from selection alone.

\figw{paths}{\linewidth}{Paths in the hierarchy morphospace: E8 over commitment at 80 agents, one line per noise level
(left), and E5 over commitment with a source moving every 100 rounds, one line per $\kappa$ (right). Hollow markers start
each path.}{paths}

\subsection{Shape is not function}
Under a moving source, runs below the lock-in edge sit at the same coordinates as true hierarchies, and cost the same: a
hub on a stale agent is as tree-like as a hub on the source. Position in the morphospace must be read together with
tracking.

\section{Relation to the original work}\label{sec:origins}

\paragraph{Source-to-structure networks.} \citet{bromberg2020hierarchy} asked how hierarchy arises from local adaptation
to where information appears. Its premise --- that the location of information shapes who listens to whom, producing
hierarchy --- holds, within a window of commitment $\kappa<T<\Thalf$. Its hard selection cut imposes hierarchy even when
the signal carries no information (with a random source the cut still climbs toward the same region of the morphospace),
and its hierarchy measure needed two corrections (\cref{app:errata}).

\paragraph{Optimal learning in noisy niches.} \citet{lee2024constructing} predicted that optimal memory grows like
$\tauE^{1/2}$ and that agents who stabilize their environment create a shared good that can favour monopolies. The
attention network reproduces the scaling once signals are noisy (exponent 0.43); with memory and pace coupled, it settles
at a stable point with free-riding; but paying for attention does not entrench hubs. Their agents correspond to the
evidence rule of \cref{sec:learning}; the network adds relays, which help under noise.

\paragraph{Adaptive transport networks.} Fluctuating loads keep loops
\citep{corson2010fluctuations,katifori2010damage}. In the attention model this holds only for a flow-based reward; with
closeness, a fluctuating source merely delays the transition to some tree. The kind of budget, not the reward, decides
robustness and cost, and transport dissipation prices the structure that the attention model builds without paying for it.

\paragraph{The synthesis note's predictions.} Three regimes (distributed, hierarchy, over-betting): confirmed, with exact
boundaries. Memory scaling: confirmed with noise. Paid-attention monopolies: not found. The update interval as a sampling
window: holds only when rewards are averaged rather than distances summed.

\paragraph{Related learning dynamics.} The update belongs to a family studied in game theory and statistical physics:
replicator dynamics with memory loss \citep{sato2003coupled}, experience-weighted attraction
\citep{camerer1999experience}, reinforcement-driven network formation
\citep{skyrms2000dynamic,bala2000noncooperative}, and adaptive networks more broadly \citep{gross2008adaptive}.

\section{Implications and limits}\label{sec:implications}

The model is abstract, and speaks to any group that must decide whom to listen to while the source of useful information
moves: teams and organizations, communities of experts, sensor and communication networks, animal groups. The following
are implications of the model; how far they carry to any real system is untested.

\begin{enumerate}
  \item \textbf{Set commitment relative to self-reinforcement, not relative to the signal.} The safe setting is just above
    the lock-in edge. Below it early advantage decides the leader whatever the evidence; the remedies are lowering
    self-reinforcement or occasionally resetting attention, not better signals.
  \item \textbf{Larger groups gain more per unit of structure but have less room for error}: coordination has increasing
    returns to size, and the tolerance for miscalibrated commitment shrinks like $1/\ln N$.
  \item \textbf{Against noisy signals, invest in signal quality or memory rather than structure.} Noise does not raise the
    cost of the hub; it lowers what the hub delivers and narrows the safe range.
  \item \textbf{Judge coordination by function, not shape.} A clear hub looks the same whether it is current or stale;
    tracking and the calibration gap tell them apart.
  \item \textbf{Side loops under noise are part of how the system copes}; removing them makes the hub more fragile.
  \item \textbf{Learning through intermediaries is cheap error correction when memory is short}; with long memory,
    careful individual inference is the benchmark.
  \item \textbf{Memory and pace co-adapt with a free-rider problem}: the stable memory is shorter than the group optimum.
\end{enumerate}

\paragraph{Limits.} Agents can listen to everyone (complete graphs) except in E6. The cost and noise studies fix
$\kappa=0.2$ and a fast memory ($\varepsilon=1$), and have not been repeated across $\kappa$. Groups have at most 80 agents
with 2--4 repeats per setting; effects reported are large compared with the spread between repeats, but differences of a
few percent are not resolved. Agents follow fixed rules. The structural cost is a proxy from transport networks, not
something the agents pay.

\paragraph{Open questions.} How much noise robustness survives on sparse graphs, where relays matter more? Would groups find
the narrow window by themselves if commitment and memory could evolve together? A model in which agents pay the
structural cost would close the loop between attention and transport. A two-class reduction (source, recent rivals, the
rest) should capture the remaining part of the noise robustness analytically.


\bibliographystyle{plainnat}
\bibliography{references}

\appendix
\section{Methods and reproducibility}\label{app:methods}

Every study follows one pipeline (\cref{fig:pipeline}): choose a grid of settings; run every cell from random starting
attention with a fresh environment draw, batched on a GPU (JAX, \texttt{jit} and \texttt{vmap} over runs,
\texttt{lax.scan} over time); record tracking, log-loss, information cost and attention snapshots every 300 rounds and
average the second half; compare with exact formulas, reduced models and switched-off variants.

\figw{pipeline}{\linewidth}{How each study runs.}{pipeline}

\paragraph{Code and data.} The Python package \texttt{nichenet} implements the model (faithful and clean modes), the
niche-learning model, transport networks, the hierarchy coordinates and the reductions; every closed form and reduction
has a test. Each experiment is a module \texttt{nichenet.experiments.<id>} that writes a JSON summary to
\texttt{data/site/}; the tables of this paper are generated from those files (\texttt{paper/tables.py}) and its figures
are exported from the companion website, which renders the same files (\texttt{site/scripts/export-figures.mjs}).
\texttt{make paper} rebuilds both.

\begin{table}[htbp]\centering\small
\caption{Experiments. All use clean mode, complete graphs (except E6), tempered selection, $\delta=10$, $\rho=2.6$.}
\begin{tabular}{lp{5.2cm}p{5.6cm}r}\toprule
id & varied & fixed & runs\\\midrule
E1 & memory ($\varepsilon$ or $\kappa$), $\tauE\in[10,3000]$ & $N=30$, $\nu=0.5$, $T=0.3$, run length $40\tauE$ & 360 per knob\\
E2 & $N\in\{10,\dots\}$, $T$, $\kappa\in\{0,0.3\}$ & constant source & 320\\
E3 & payment $a\in\{0,0.5,1,2\}$, $\tauE\in\{200,2000\}$ & $N=30$, $T=0.6$; $\kappa$ copied from best earner & 32 $\times$ 150 epochs\\
E4 & resident and invader memory, agent type, $\tauE\in\{100,1000\}$ & binary environment, invaders at 1\% & 2 seeds $\times$ $10^7$ steps\\
E5 & $T$ (14), $\kappa\in\{0,0.2,0.4,0.6\}$, source schedule (5) & $N=30$, 3000 rounds, no noise & 840\\
E6 & budget, usage, $p\in[0.4,20]$, schedule & $7\times7$ lattice, $\eta=0.05$, 3000 steps & 720\\
E7 & pooling window $k$, $\tauE$, rule & $N=30$, $\nu=0.5$ & 180 per rule\\
E8 & $N\in\{10,20,40,80\}$, $\nu\in\{0,\dots,0.75\}$, $T$ (7) & $\kappa=0.2$, $\varepsilon=1$, $\tauE=300$, 3000 rounds & 224\\
E9 & reward rule (3) + Bayes filter & E8 grid, seeds and environments & 672 + 32\\
A4 & reward rule $\times$ peers free/equalized & constant source, $N$ (4) $\times$ $\nu$ (3) $\times$ $T$ (4), 4000 rounds & 384\\\bottomrule
\end{tabular}\end{table}

\section{Derivations}\label{app:derivations}

\paragraph{B2: the update as an Euler step.} Writing $R(A)$ for the reward matrix, $A'=(A+\varepsilon R)/(1+\varepsilon)
=A+\frac{\varepsilon}{1+\varepsilon}(R(A)-A)$, a forward-Euler step of $\dot A=R(A)-A$ with step
$\varepsilon/(1+\varepsilon)$.

\paragraph{B3: fixed points.} At a fixed point $A=R(A)$, i.e.\ $A_{ki}\propto(D_k^{\kappa-1}A_{ki}^\kappa)^{1/T}$. Taking
logarithms, $(1-\kappa/T)\ln A_{ki}=-\frac{1-\kappa}{T}\ln D_k+c_i$, so $A_{\cdot i}=\mathrm{softmax}(-\ln D/\Teff)$ with
$\Teff=(T-\kappa)/(1-\kappa)$. This is the rest point of the Sato--Crutchfield replicator with reward $-\ln D$
\citep{sato2003coupled}; a 600-step simulation reaches it to $10^{-9}$.

\paragraph{A1: exact symmetric reduction.} With a constant source $s$ on a complete graph and the symmetric state (every
$i\ne s$ puts $x$ on $s$ and $p=(1-x)/(N-2)$ on each other agent; $s$ spreads its attention evenly), row $s$ of $P$ is
uniform over $k\ne s$. The permutations of $\{1,\dots,N\}\setminus\{s\}$ map the state to itself, so $D_{sk}=D$ for all
$k\ne s$ at any $\delta$, and $D_{ss}=D/\rho$. The reward ratio of source to peer is
$\big[(D/\rho)^{\kappa-1}x^\kappa/(D^{\kappa-1}p^\kappa)\big]^{1/T}=(\rho^{1-\kappa}u^\kappa)^{1/T}$, independent of $D$.
The reduced map is $x\mapsto\big(x+\varepsilon\,a/(a+(N-2)b)\big)/(1+\varepsilon)$ with $a=(\rho^{1-\kappa}x^\kappa)^{1/T}$,
$b=(p^\kappa)^{1/T}$; its fixed point is $u^*=\rho^{(1-\kappa)/(T-\kappa)}$ and its slope there
$(1+\varepsilon\kappa/T)/(1+\varepsilon)$. Setting $x^*=1/2$ gives $u^*=N-2$ and hence $\Thalf$.

\paragraph{A2: stability through the kink.} Because $D_{ss}=\min_kD_{sk}/\rho$ and all $D_{sk}$ tie at the symmetric
state, the map is piecewise smooth with $N-1$ pieces, each anchored at a relay $k^*$ attaining the minimum; all have the
same spectrum by symmetry. We linearize the piece anchored at relay 1 in the tangent space of the per-agent budgets
($N(N-2)$ directions) by central finite differences scaled to the smallest weight. The collective mode contracts at
exactly $(1+\varepsilon\kappa/T)/(1+\varepsilon)$; most modes sit on this bulk value and a relay group lies 0.01--4\%
above it (\cref{tab:spectra}). A finite-difference Jacobian taken across the kink mixes pieces and produces a spurious
eigenvalue (0.75 at $N=10$, $\kappa=0.3$, $T=0.6$).

\paragraph{A3: separatrix.} For two dominant entries $w_s,w_q$ of a column, the reward ratio is
$(\rho^{1-\kappa}(w_s/w_q)^\kappa)^{1/T}$, so $l=\ln(w_s/w_q)$ evolves toward $\frac{1-\kappa}{T}\ln\rho+\frac{\kappa}{T}l$
under the reward; its fixed point $l^*=(1-\kappa)\ln\rho/(T-\kappa)$ is repelling for $T<\kappa$. Lock-in onto $q$ happens
iff $w_q/w_s>\rho^{(1-\kappa)/(\kappa-T)}$ initially; as $T\to0$ this becomes $(w_q/w_s)^\kappa>\rho^{1-\kappa}$, the
one-step criterion of the hard cut. Columns interact weakly through relay distances, which is why slightly fewer
listeners lock in than predicted at the lowest $T$ (\cref{tab:separatrix}).

\paragraph{A4: noisy mean field and its refinements.} With noise, the true source is seeded with probability
$1-\nu+\nu/N$; otherwise a peer $j$ carries the $\rho$ advantage and $s$ competes as a strong peer:
\begin{align}
  f(x)&=\frac{a}{a+(N-2)b},\qquad g(x)=\frac{c}{\rho^{(1-\kappa)/T}b+c+(N-3)b},\qquad c=x^{\kappa/T}.
\end{align}
Three refinements: (i) the averaged map $x=\mathbb E[r](x)$; (ii) the stochastic walk with the noise drawn each round
(fluctuations raise the time average by Jensen's inequality); (iii) the relay-aware walk, in which the distances in a
noise round are measured from $j$: $s$'s own attention is spread ($q=1/(N-1)$), $j$'s audience share toward $s$ is
$q/(q+(N-2)p)$, the largest, so $D_s=\delta-\ln[q/(q+(N-2)p)]$ is the shortest relay distance and $s$ keeps more reward.
The remaining gap is closed by peer heterogeneity, tested in \cref{sec:robustness}: the column walk keeps one agent's
full unequal column, and its competition factor $\sum_kw_k^{\kappa/T}/[(N-2)\bar w^{\kappa/T}]<1$ measures the concavity
effect.

\section{Validation and errata}\label{app:validation}

\paragraph{Validation.} Before any new experiment the components were checked against the original work:
\begin{itemize}
  \item The faithful mode equals the original code step for step on 28 cases (attention matrices after every update and
    effective distances, relative tolerance $10^{-9}$); the original hierarchy and efficiency coordinates agree on test
    matrices.
  \item The niche-learning eigenfunction solver equals the authors' on six cases (relative tolerance $10^{-8}$); the
    scaling $\taum^*\sim\tauE^{1/2}$ of their Fig.~2 is reproduced; their invasion results are reproduced for monopolized
    niches and partly for shared niches.
  \item Transport networks: steady loads give spanning trees for $\gamma<1$ and keep all edges for $\gamma>1$;
    fluctuations create loops, rising with $\gamma$; adaptation reaches Corson's optimum.
  \item The GPU kernel equals the reference implementation on seven configurations; closed forms B2, B3, B5 and A1--A4 are
    tests.
  \item One thesis claim was not reproduced as stated: on complete graphs the selection transition does not sit at a fixed
    number of kept entries.
\end{itemize}

\paragraph{Errata in the original code}\label{app:errata}
\begin{description}
  \item[H1. Element-wise instead of matrix powers.] The graph entropies of the hierarchy coordinates sum walk
    probabilities $P=\sum_k(B^\top)^k$, a sum of matrix powers (eq.~14 of \citealp{corominas2013originsheirarchy}). The code
    used element-wise powers, whose $k=0$ term is the all-ones matrix, so the entropies reduce to plain sums of
    log-degrees.
  \item[H2. Treeness over the wrong set of subgraphs.] Treeness averages over the network once plus the pruned subgraphs
    that still contain a link; the code counted the network twice and included empty subgraphs.
  \item[Effect.] On 180 evolved networks the code underestimates treeness by 0.1--0.25 on average (H2 dominates);
    under a constant source corrected treeness is near 0.7 for every selectivity, where the code showed a flat 0.45 with a
    jump at full selectivity. Feedforwardness and orderability are unaffected.
\end{description}
Implementation behaviours of the original code that change results without contradicting a published definition
(rewards on non-edges at $\kappa=0$, ranking over all entries rather than in-neighbours, mixed normalizations, fixed
$\varepsilon=1$, a source index that silently meant random seeding, an unseeded random generator) are kept in faithful mode
and fixed in clean mode.

\section{Glossary}\label{app:glossary}
\input{generated/glossary}

\end{document}
